\documentclass[11pt]{amsart}

% ── Packages ──────────────────────────────────────────────────────────
\usepackage{amsmath,amssymb,amsthm}
\usepackage[utf8]{inputenc}
\usepackage[T1]{fontenc}
\usepackage{hyperref}
\usepackage{booktabs}
\usepackage{enumitem}
% \usepackage{microtype}  % requires scalable fonts; enable if using lmodern

\hypersetup{
  colorlinks=true,
  linkcolor=blue,
  citecolor=blue,
  urlcolor=blue,
}

% ── Metadata ──────────────────────────────────────────────────────────
\title{Substrate Scanning Engine:\\
  Geometric Parameter Generation for $L$-functions\\
  via Elementary Visualization and Systematic Elimination}

\author{Ricardo Hern\'andez Reveles}
\address{Independent researcher, Guadalajara, M\'exico}
\email{} % TODO: add if desired
\urladdr{https://orcid.org/0009-0001-1993-1471}
\date{March 2026}
\thanks{Serie~I Working Paper. Code: AGPL-3.0-or-later. Papers: CC-BY-4.0.}

\subjclass[2020]{11M06 (primary); 68T05, 65D99 (secondary)}
\keywords{Riemann zeta function, spectral features, geometric machine learning,
  Euler product, $L$-functions, Substrate Scanning, computational phenomenology}

\begin{document}

\begin{abstract}
We describe a modular computational engine that takes any Euler product and
produces geometric parameters for machine learning without statistical
assumptions. The engine maps the Euler product to a toroidal embedding,
decomposes the resulting dynamics into three separable scales (drift, spiral,
continued-fraction modulation), constructs the electrostatic potential
$V(\sigma,t) = -\log\lvert -\zeta'/\zeta \rvert$, computes its Laplacian
(confirming harmonicity between zeros at $0.07\%$ median residual), and
extracts 29 deterministic features per query point. Applied to the Riemann
zeta function, the engine discriminates zeros from non-zeros with $5.6\sigma$
separation on phase coherence alone, using only the partial Euler product.

The central decomposition is
$\theta^*(k,\gamma) = f(b)\cdot\cos(\ln(p)\cdot\gamma)$, where
$f(b) = b/(b+1)$ is the natural per-step contraction (universal,
amplitude-only) and $\cos(\ln(p)\cdot\gamma)$ is the phase modulation
(specific to each zero). Three independent tests confirm that the arithmetic
channel (factorization of $p-1$) is orthogonal to both amplitude and phase
channels, establishing that shared algebra does not imply shared information.

The engine is demonstrated on 16 domains: the Riemann zeta function,
Dirichlet $L$-functions, elliptic curve $L$-functions, Cayley graph spectra,
4-port differential S-parameters (signal integrity), vibration spectra
(rotating machinery fault diagnosis), 2D physical cavities
(antenna/acoustic/laser optimization), visible-light spectroscopy
(colorimetry), general graph analysis (social networks, supply chains, bot
detection), abstract syntax trees (code clone detection, refactoring), phonon
spectra (materials design, thermoelectric screening), 3D wavelet subband
analysis (volumetric data), turbulence cascade classification, 2D image
spectral analysis (forgery detection), spectral pathfinding (graph
navigation), and elliptic curve cryptographic auditing---with zero
architectural changes between domains. The constraint is precise: the system
must have a linear operator with discrete spectrum. Under this condition, the
motor generates, ML selects, and domain physics verifies. Estimated
computational savings range from $10\times$ (BSD verification) to
$10{,}000\times$ (materials design: 29,241 compositions screened in 32
seconds vs.\ DFT).

No result in this paper is claimed as mathematically novel. The contribution
is methodological: a systematic pipeline for geometric exploration that
produces falsifiable outputs, documents its own overclaims, and eliminates
impossible approaches at visual speed.
\end{abstract}

\maketitle


% ══════════════════════════════════════════════════════════════════════
\section*{0.\ Errata and Methodological Transparency}
\label{sec:errata}
\addcontentsline{toc}{section}{0.\ Errata and Methodological Transparency}
% ══════════════════════════════════════════════════════════════════════

This section appears first because all subsequent claims should be read in
light of these corrections. During the development of this engine, seven
overclaims were produced and corrected within the same working sessions. Each
is documented below with the quantitative evidence that falsified it.

\subsection*{0.1 Overclaims Corrected}

\begin{enumerate}[label=(\alph*)]
\item \textbf{``Exactly as $f(b)$''} --- the per-step contraction
  $f(b) = b/(b+1)$ is algebraically exact, but the cumulative exponent is
  $-1.4 \pm 0.1$, not $-1.0$. The discrepancy arises because per-step
  exactness does not imply exact geometric series accumulation.

\item \textbf{Spiral exponent $\alpha = -0.96$} --- artifact of
  fixed-reference measurement ($P(N_{200})$ as anchor).
  Reference-independent methods (moving center, consecutive differences) give
  $\alpha = -1.4 \pm 0.1$ consistently across 5 methods and 8~zeros.

\item \textbf{Poincar\'e compactification produces convergence to
  $\partial B^3$} --- incorrect. Since $|P(N)|$ decreases with~$N$ (the
  normalizing sum $Z(N)$ grows in the denominator), the compactified points
  move toward the origin of~$B^3$, not toward the boundary. The gap
  $1-|\tilde{U}_P|$ grows from $0.59$ ($N=10$) to $0.70$ ($N=2000$). The
  algebraic identity $f(b) = \text{Poincar\'e}_{1D}$ remains valid; the
  geometric application to these trajectories was the overclaim.

\item \textbf{The $f(b)$ bridge between digit-sum and toroidal channels} ---
  empirically falsified by the dual test ($v_2(p-1)$ as cross-channel
  predictor, 5132 primes): $v_2(p-1)$ predicts digit-sum deficit with
  $\rho = -0.43$ but has $\rho < 0.01$ for all toroidal quantities.

\item \textbf{$\gamma \cdot |\omega| \approx \text{constant}$} ---
  $\text{CV} = 36\%$, not constant. Removed.

\item \textbf{The cone of radial directions in $B^3$ closes with~$N$} ---
  it does not. Semi-angle decays as $N^{-0.19}$, too slow for convergence.

\item \textbf{``The zeros exist as points on $\partial B^3$''} ---
  reification of the embedding. The zeros exist as zeros of the analytic
  continuation of~$\zeta(s)$. The compactified trajectories point toward
  boundary positions that encode the zeros; they do not constitute the zeros.
\end{enumerate}

The ratio of eliminated claims to surviving results is $7{:}1$. This ratio is
characteristic of honest computational phenomenology and is the primary
diagnostic of the methodology.


% ══════════════════════════════════════════════════════════════════════
\section{Introduction: $t$ as Parameter}
\label{sec:intro}
% ══════════════════════════════════════════════════════════════════════

The factorization
\begin{equation}\label{eq:fundamental}
  p^{-s} = p^{-\sigma} \cdot e^{-i\ln(p)\cdot t}
\end{equation}
is a fact, not a choice. The parameter~$t$ traces the critical line; each
prime generates a mode with frequency $\omega = \ln(p)$; the zeros of the
Riemann zeta function are the values of~$t$ where all phases conspire to
annihilate the sum. Everything else in this paper---the embedding, the three
scales, the potential landscape, the envelope decomposition---is scaffolding
around this single equation.

The Euler product $\prod_p (1 - p^{-s})^{-1}$ converges absolutely for
$\sigma > 1$. At $\sigma = 1/2$, the partial products do not converge
(Conrad~\cite{Conrad2005}, extending earlier work of Goldfeld and Kurokawa
et~al.). The non-convergence is not a technical obstacle---it is the
fundamental structural fact. Kurokawa's ``Deep Riemann Hypothesis'' asserts
that convergence of partial Euler products at the center of the critical
strip would imply something strictly deeper than the generalized Riemann
hypothesis.

This paper describes a computational engine that takes any Euler product and
produces its geometric atlas: a collection of deterministic features derived
from the product's structure, suitable as input for machine learning without
statistical assumptions.


% ══════════════════════════════════════════════════════════════════════
\section{The Geometric Engine}
\label{sec:engine}
% ══════════════════════════════════════════════════════════════════════

The engine consists of eight modules, each independently testable. The input
is an Euler product (abstract base class); the output is a 29-dimensional
feature vector per query point $(\sigma, t, K)$.

\subsection{Module 1: Euler Product}

Abstract interface with methods: \texttt{primes(K)},
\texttt{local\_factor(p,s)}, \texttt{log\_deriv\_term(p,s)},
\texttt{arithmetic\_channel(p)}. Concrete implementations provided for
$\zeta(s)$, $L(s,\chi)$, and $L(s,E)$. The \texttt{arithmetic\_channel}
method returns features specific to each $L$-function: $v_2(p-1)$ and
$\mathrm{ord}_2(p)/(p-1)$ for~$\zeta$; $\chi(p)$ for Dirichlet; $a_p$ for
elliptic curves.

\subsection{Module 2: Toroidal Embedding}

Maps the Euler product to weighted coordinates on the mixed torus
$\mathbb{T}^2_{(p_k, p_{k+1})}$ via $p^{-s} = p^{-\sigma}\cdot e^{-i\omega
t}$. The embedding is standard (Bohr torus, 1920s). What is specific to this
work is the pair-weighted parametrization and the decomposition into three
scales.

\subsection{Module 3: Three Scales}

The trajectory $P(N)$ decomposes into: Scale~1 (drift, $\sim 1/\ln N$,
divergent by Mertens), Scale~2 (spiral, $\sim N^{-1.4}$,
reference-independent), Scale~3 (CF modulation, classical Khinchin). The
decomposition is specific to this embedding; the scales may not have
counterparts in other parametrizations.

\subsection{Module 4: Envelope Decomposition}

The natural per-step weight is $f(b_k) = b_k/(1+b_k)$ where
$b_k = p_k^{-\sigma}/Z_k$. Optimization of weighted log-derivative loss over
known zeros produces optimal weights $\theta^*(k,\gamma)$. The central result
of the pipeline:
\begin{equation}\label{eq:envelope}
  \theta^*(k, \gamma)
  = f(b_k)\cdot\bigl[A + B\cos(\ln(p_k)\cdot\gamma)
                         + C\sin(\ln(p_k)\cdot\gamma)\bigr]
\end{equation}
where $f(b)$ is universal (amplitude decay) and the trigonometric factor is
zero-specific (phase tuning). $R^2$ ranges from $0.09$ ($K=20$) to $0.31$
($K=30$, $\gamma_3$). The phase factor is the numerical fingerprint of
Connes' prolate operator.

\subsection{Module 5: Potential Landscape}

$V(\sigma,t) = -\log\lvert -\zeta'/\zeta(\sigma+it)\rvert$ is the
electrostatic potential with charges at the zeros. Between zeros,
$\Delta V \approx 0$ (median $|V|/|V| = 0.0007$ over 458 grid points). Near
zeros, $\Delta V$ diverges. The integral $\oint V\,dA$ around each zero
approaches $2\pi$ in principle (Dirac delta charges), though for the partial
product the charges are diffuse rather than point-like.

\subsection{Module 6: Weil Spectrum}

Discretized Weil quadratic form
$W_{jk} = \sqrt{\ln p_j}\cdot\sqrt{\ln p_k}\cdot p_j^{-1/4}\cdot
p_k^{-1/4}\cdot\cos\bigl((\ln p_j - \ln p_k)\cdot\gamma\bigr)$.
Eigenvalues computed by Jacobi iteration for $K \le 150$. The minimum
eigenvalue detects the zero; the spectral gap measures resolution.

\subsection{Module 7: Channel Orthogonality}

Tests independence between Canal~A (arithmetic: $v_2$, $\mathrm{ord}_2$,
$\chi$, $a_p$), Canal~B (amplitude: $\ln p$, $p^{-1/2}$), and Canal~C
(phase: $\cos(\ln p \cdot \gamma)$, $\sin(\ln p \cdot \gamma)$). Maximum
Spearman~$\rho$ between column pairs. For~$\zeta$: $A\perp B$
($\rho < 0.28$), $A\perp C$ ($\rho < 0.33$), $B$ not orthogonal to~$C$
(share $\ln p$).

\subsection{Module 8: Feature Vector}

Combines all modules into a 29-dimensional vector: 3~meta ($\sigma$, $t$,
$K$), 4~torus, 6~potential ($V$, $\nabla V$, $\Delta V$, harmonicity),
5~Weil (eigenvalues), 6~envelope ($f(b)$ stats, phase coherence),
2~Poincar\'e (radius, gap), 3~channel ($\rho$ cross-correlations). Output is
a \texttt{numpy} array, deterministic, reproducible, and computable in
${\sim}8\,$ms per query.


% ══════════════════════════════════════════════════════════════════════
\section{Results}
\label{sec:results}
% ══════════════════════════════════════════════════════════════════════

\subsection{Zero Discrimination}

At $K=30$, $\sigma=0.5$, the three features with highest zero/non-zero
separation are: phase coherence ($5.6\sigma$), $V$~potential ($2.4\sigma$),
and gradient magnitude ($0.75\sigma$). Phase coherence alone classifies zeros
with zero training. This is the Canal~C signal: at a zero, the phases
$\ln(p)\cdot\gamma$ conspire, producing high coherence; away from zeros, they
cancel.

\subsection{Envelope Factorization}

Per-zero optimal weights are \textsc{zero-specific} (mean pairwise
$|r| = 0.14$). The third axis beyond Canal~A and~B is
$\cos(\ln(p)\cdot\gamma)$---the phase of the target zero evaluated at each
prime ($\rho = -0.59$, $p = 6.6\times 10^{-4}$ at $K=30$). Canal~A ($v_2$,
$\mathrm{ord}_2$) is absent from the optimal weights at every~$K$ tested
($r < 0.1$). The optimal operator for zero detection is $f(b)\cdot\cos(\omega)$,
which is purely spectral.

\subsection{The $\sigma$ Transition}

Animating $\sigma$ from 1.2 to 0.3 shows the convergence/divergence
transition as a continuous geometric deformation: the spiral passes from
closed (convergent product) to open (divergent drift). The transition point
is the boundary of the Euler product's convergence region, visible as the
moment the spiral stops closing. This is the prime number theorem made curve:
the spiral does not close at $\sigma = 1/2$ because $\sum 1/p$ diverges
(Mertens), and the pseudo-randomness of primes decreases with~$p$, visible
as increasing coherence of consecutive steps.

\subsection{Geodesics}

Gradient descent on $V(\sigma,t)$ from any $\sigma > 1/2$ converges to
$\sigma = 1/2$. The geodesics channel between zero-wells, attracted by the
saddle points. The surface~$V$ is harmonic between singularities
($\Delta V = 0$), confirming the electrostatic analogy. The zeros are point
charges; the critical line is the equilibrium of the boundary conditions
imposed by the functional equation's symmetry $s \leftrightarrow 1-s$.

\subsection{Cross-Domain Validation}

The engine runs on elliptic curve $L$-functions with zero code changes. For
four toy curves ($y^2 = x^3 + ax + b$), the potential~$V$ at $\sigma=1$
discriminates rank: the curve with a rational point (suggesting rank
$\ge 1$) has $V = -1.32$, deeper than rank-0 curves ($V \approx -1.00$ to
$-1.23$). Canal~A ($a_p$) is orthogonal to all geometric features
($\rho < 0.11$). The engine also maps to Cayley graph spectra: eigenvalues
of the graph Laplacian correspond to zeta zeros, the spectral gap
corresponds to the zero-free region, and the Ramanujan bound corresponds to
RH.


% ══════════════════════════════════════════════════════════════════════
\section{Connection to Recent Work}
\label{sec:connection}
% ══════════════════════════════════════════════════════════════════════

\subsection{Connes (arXiv:2602.04022, February 2026)}

Connes~\cite{Connes2026} extremizes the Weil quadratic form restricted to primes $\le 13$,
approximating the first 50~zeros with precision ranging from
$2.6\times 10^{-55}$ to~$10^{-3}$. His prolate operator is the
formalization of our Module~4: it constructs optimal weights by extremizing a
quadratic form. Our decomposition $\theta^* = f(b)\cdot\cos(\omega)$ is the
numerical separation of his operator into infrastructure ($f(b)$, amplitude)
and signal ($\cos(\omega)$, phase). Our engine achieves $10^{-6}$ precision
with L-BFGS on the log-derivative loss; the 49 orders of magnitude
difference reflects the difference between the Weil form and a naive loss
function.

\subsection{Conrad (\textit{Canadian J.\ Math.}, 2005)}

Conrad proved that convergence of partial Euler products on the critical line
is equivalent to a condition strictly deeper than the Riemann hypothesis. Our
computational verification (50-digit Cauchy test, three independent methods)
confirms his theorem. Our geometric contribution is making the
non-convergence visible as a spiral that never closes.

\subsection{Guth--Maynard (2024, accepted \textit{Annals} 2025)}

Their zero density estimate~\cite{GuthMaynard2024}
$N(\sigma,T) \le T^{30(1-\sigma)/13 + o(1)}$ improves Ingham's 1940 bound.
The result has no direct intersection with our work: their techniques (large
value estimates, decoupling) operate in a different space from our geometric
features. We include it for completeness, not for borrowed relevance.


% ══════════════════════════════════════════════════════════════════════
\section{Attribution and Scope}
\label{sec:attribution}
% ══════════════════════════════════════════════════════════════════════

\begin{enumerate}[label=(\alph*)]
\item The non-convergence of partial Euler products at $\sigma = 1/2$ is
  Conrad~\cite{Conrad2005}, formalizing earlier work. Our contribution:
  computational verification and geometric visualization.

\item The toroidal embedding via $p^{-s}$ on the Bohr torus is classical
  (Bohr 1920s). Our contribution: the pair-weighted mixed-torus
  parametrization and three-scale decomposition.

\item The identity $f(b) = b/(b+1)$ as the Poincar\'e compactification in
  1D is elementary. Our contribution: identifying it as the per-step
  contraction of the Euler product, and proving that its geometric
  application fails (gap grows, cone does not close).

\item The class number identity $S - k/2 = -h(-p)/2$ is derivable from
  Dirichlet 1839. Our contribution: the route of discovery (Substrate
  Scanning through 5~phases) and explicit verification across 961~primes.

\item The electrostatic interpretation of $\log\lvert -\zeta'/\zeta \rvert$
  is implicit in the theory of meromorphic functions. Our contribution:
  quantitative verification of harmonicity ($0.07\%$ median) and the
  visualization as 3D surface with Laplacian coloring.
\end{enumerate}

No result in this paper is claimed as a new theorem. The contribution is the
engine and the methodology that produced it.


% ══════════════════════════════════════════════════════════════════════
\section{Methodology: Substrate Scanning}
\label{sec:methodology}
% ══════════════════════════════════════════════════════════════════════

Substrate Scanning is rapid adversarial computational auditing that produces
strategic intelligence by systematically eliminating approaches that cannot
work. The cycle is: observation $\to$ overclaim $\to$ quantitative
verification $\to$ correction $\to$ audited result.

In this paper, the cycle executed 7~times. Each overclaim was produced by the
same failure mode: a partial correspondence was generalized before
verification of the boundary conditions. Each correction was triggered by a
specific quantitative test: doubling~$N$ (killed $P(\infty)$), computing
reference-independent exponents (killed $\alpha = -0.96$), the dual test
(killed the channel bridge), the cone analysis (killed Poincar\'e
convergence).

The methodology's value is not in producing results but in the speed of
elimination. Seven approaches that might take months to evaluate analytically
were eliminated in a single computational session, each with quantitative
evidence. The engine is the distillation of what survived.


% ══════════════════════════════════════════════════════════════════════
\section{Applications}
\label{sec:applications}
% ══════════════════════════════════════════════════════════════════════

\subsection{BSD Verification}

The engine generates 29 geometric features per $(\sigma, t)$ query for any
elliptic curve $L$-function. A classifier trained on feature vectors near
$s = 1$ for curves of known rank could predict rank for curves of unknown
rank. The test is binary and falsifiable: does the $V(\sigma)$ profile near
$s = 1$ discriminate rank~0 from rank $\ge 1$? Toy results (4~curves,
$K=20$) show $V = -1.32$ for the rank-1 candidate vs.\
$V \approx -1.00$ for rank-0 curves.

\subsection{Optimal Prime Selection}

The decomposition $\theta^* = f(b)\cdot\cos(\ln(p)\cdot\gamma)$ identifies
which primes contribute most to detecting a specific zero. For numerical
computation of $L$-function zeros, this reduces cost: 6 well-chosen primes
give what 50 generic primes cannot (Connes' result as practical heuristic).

\subsection{Cayley Graph Classification}

The engine maps directly to Cayley graph spectra. Eigenvalues of the graph
Laplacian are computed; the spectral gap is measured; the Ramanujan bound is
checked. A classifier trained on geometric features could detect expander
graphs without full diagonalization---useful in cryptography and network
design.

\subsection{Signal Integrity: S-Parameter Classification}

The S-matrix $S(f)$ of a 4-port differential channel is a spectrum
parametrized by frequency~$f$. The mixed-mode decomposition into $S_{dd}$
(differential), $S_{cc}$ (common mode), and $S_{dc}$ (mode conversion) maps
exactly to Canal~B, Canal~A, and cross-channel respectively. The motor
generates 29~features from $S(f)$ in $0.58\,$ms per file, classifying
channels as PASS/MARGINAL/FAIL at 1709 classifications/second. Validated on
5~synthetic profiles: resonant channels, lossy channels, and high-crosstalk
channels are correctly identified with diagnostic reasons. The channel
orthogonality test quantifies mode conversion coupling---the signal integrity
equivalent of ``shared algebra does not imply shared information.'' For
production SI labs measuring 500 boards/day, this replaces manual
S-parameter curve inspection with automated classification.

\subsection{Rotating Machinery Fault Diagnosis}

Vibration spectra $V(f)$ from accelerometers on rotating machinery have
harmonic orders $n\times$ (analogous to primes), fault frequencies
BPFO/BPFI/BSF (analogous to zeros of~$\zeta$), and an envelope of harmonic
decay (analogous to $f(b)$). The sideband modulation pattern around fault
frequencies is $\cos(\omega)$---the phase tuning. The motor generates
29~features from the vibration spectrum at 1132 diagnoses/second ($0.88\,$ms
per sample), classifying: normal, unbalance, misalignment, mechanical
looseness, bearing outer/inner race defects, and gear mesh faults. Real-time
at 50\,kHz sampling with 8192-point FFT windows ($0.16\,$s). Severity
progression is tracked: a developing bearing defect is detected consistently
from $\text{SNR}=60\,$dB through $\text{SNR}=20\,$dB as the number of
non-synchronous peaks grows.

\subsection{Physical Cavity Optimization}

The eigenvalues of the Laplacian on a 2D cavity are the resonant
frequencies---directly analogous to zeros of~$\zeta$. Weyl's law is the
analogue of the zero density estimates. The spectral gap determines the
quality factor (antenna), the evenness of frequency response (acoustics), or
the mode isolation (laser). Tested on 8~cavity shapes (circle, ellipse,
rectangle, stadium, parabolic, L-shape): the motor correctly ranks parabolic
as optimal for antenna (high~$Q$), stadium for acoustic rooms (uniform
spacing), ellipse for laser cavities (large gap), and stadium for quantum
billiards (high level repulsion, confirming the Bunimovich stadium as the
classical chaotic cavity). Cross-domain comparison shows the ellipse has
level spacing ratio $\langle r\rangle = 0.53$, closest to GUE ($0.536$) and
to $\zeta$ zeros ($0.581$)---a numerical echo of the Montgomery--Odlyzko
conjecture. Estimated computational savings: factor ${\sim}150\times$ versus
full diagonalization when sweeping $10^4$ candidate geometries.

\subsection{Colorimetry and Spectral Compression}

A visible spectrum $S(\lambda)$ sampled at 41~points (380--780\,nm, 10\,nm
intervals) is compressed to 3~values by CIE~XYZ, losing
metamer-distinguishing information. The engine compresses to 29~features
preserving spectral structure the eye discards: peak positions, bandwidth,
spectral entropy, purity, and an effective metamerism dimension. Validated on
11~synthetic spectra (pure R/G/B, daylight, incandescent, fluorescent, LED
white, ruby, emerald, and two designed metamer pairs). The metamer pair test
confirms that feature distance ($52.8$) exceeds XYZ distance ($26.6$) by a
factor of $2.0\times$---the features detect differences that XYZ discards.
Throughput: 1677 spectra/second ($0.6\,$ms). Standalone engine:
\texttt{colorimetry\_engine.py}.

\subsection{Wavelet 3D Representations and World Models}

Wavelets 3D decompose volumetric data into 8~subbands per level, separating
surfaces (curvelets), filaments (ridgelets), and points (isotropic wavelets)
at each scale. The motor generates 29~features per subband, producing a
tensor of geometric descriptors for the entire volume. Current state of the
art (WaveFormer, MICCAI 2025~\cite{WaveFormer2025}; WDM,
2024~\cite{WDM2024}) uses wavelets 3D as compression layers inside neural
networks---the wavelet is a computational tool, not a representation. The
gap: no existing work uses wavelets 3D as a deterministic geometric feature
generator for world models. The motor fills this gap by producing
interpretable, deterministic spectral coordinates from volumetric data before
any learning occurs. Applications include: 3D industrial CT (defect
classification by topological type---pore vs.\ crack vs.\ inclusion),
turbulence diagnosis via wavelet scale decomposition (the Kolmogorov cascade
is the ``Euler product'' of the flow), seismic interpretation (faults,
channels, and domes as geometric primitives), and 4D flow MRI (vortex
detection as ``zeros'' of the vorticity field). The motor's constraint---the
system must have a linear operator with discrete spectrum---is satisfied by
wavelets, which linearize multi-scale structure by construction.

\subsection{General Graph Analysis}

The Laplacian $L = D - A$ of any graph has discrete spectrum. The engine
computes 29~features including algebraic connectivity ($\lambda_1$), spectral
gap, level repulsion, partition balance via the Fiedler vector, effective
resistance, and hub dominance. Validated on 9~graph configurations:
Erd\H{o}s--R\'enyi sparse and dense (correctly classified as EXPANDER at
high density), Barab\'asi--Albert (SCALE-FREE, FRAGILE), Watts--Strogatz
(REGULAR), stochastic block model (MULTI-COMMUNITY), star, path, grid, and a
synthetic bot farm. Bot farm vs.\ organic network separates on 12~features
with diff $> 0.3$---the strongest being health score ($1.36$ separation),
max eigengap position ($1.11$), and small-world indicator ($0.77$). Critical
node detection via Fiedler vector identifies supply chain vulnerabilities and
network bottlenecks. Throughput: 833 graphs/second at $N=100$, $80\,$ms at
$N=500$.

\subsection{Abstract Syntax Tree Analysis}

A program's AST is a tree, a tree is a graph, and the engine applies
directly. Python source code is parsed to AST via the standard \texttt{ast}
module, converted to a labeled graph (nodes = operations, edges =
parent-child), and the 29~graph spectral features are extracted. 15
code-specific features are added: depth, branching, type entropy, control
flow ratio, and function ratio, yielding 44~features per source file.
Validated on 7~code samples ranging from clean functions to deeply nested,
spaghetti, obfuscated, and well-structured code. Clone detection correctly
identifies fibonacci/clone\_fibonacci as structurally identical (similarity
$= 1.000$). Classification detects: DEEP-NESTING (depth $= 14$), MONOLITHIC
(0 functions in 122~nodes), and WELL-MODULAR patterns. Critical AST nodes
identify refactoring targets. Current limitation: the raw spectral similarity
is too coarse for tree discrimination---all trees have similar Laplacian
shape. Normalization by size and relative eigenvalue ratios (analogous to
$\theta^*/f(b)$) is the known fix. Throughput: 139 analyses/second.

\subsection{Materials Design via Phonon Spectra}

The dynamical matrix $D(q)$ of a crystal has eigenvalues $\omega^2(q)$---the
phonon frequencies. The engine generates 29~features from the phonon spectrum
of a 1D chain model: band gap, group velocity, Debye temperature proxy,
thermal conductivity proxy, optical branch flatness (rattler modes), and
level spacing statistics. Validated on 12~materials ranging from monatomic
metals (Cu, Fe, Al) to diatomic semiconductors (NaCl, GaAs, SiC, PbTe),
polymers, perovskites, and clathrates (skutterudite proxy). The engine
correctly classifies SiC as SEMICONDUCTOR + STIFF (high gap, high
$v_\text{sound}$), PbTe and BiTe as thermoelectric candidates (low~$\kappa$,
flat optical branches), and polymer-like chains as INSULATOR (large gap from
bond alternation). Thermoelectric screening ranks materials by a $ZT$ proxy
$= (\text{gap} + 0.1)\times(\text{flatness} + 0.1)/(\kappa + 0.1)$,
correctly placing clathrates and PbTe-like materials at the top. Inverse
design: 29,241 compositions swept in 32~seconds ($905\,$/s), finding the
optimal mass/spring combination for a target spectrum. The cross-domain
mapping closes: wavevector~$q$ is~$t$, phonon frequencies are zeros, the
band gap is the zero-free region, and the Bloch phase $e^{-iqa}$ is
$\cos(\ln(p)\cdot\gamma)$.

\subsection{3D Wavelet Subband Analysis}

A 3D discrete wavelet transform (Haar, manually implemented) decomposes
volumetric data into 8~subbands per level across 3~levels. The engine
generates 29~features from subband energies, kurtosis, entropy, sparsity,
anisotropy per level, directional energy decomposition (edge/cross/corner),
scale self-similarity, and spectral centroid. Validated on 6~synthetic
volumes: sphere, cylinder, two spheres (multi-scale), Gaussian noise,
sphere$+$noise (denoising), and torus. The classifier correctly separates
STRUCTURED from NOISY (kurtosis $+$ sparsity) and identifies MULTI\_SCALE
objects. Cylinder and torus are ambiguously classified as MULTI\_SCALE
because they possess two inherent length scales; this ambiguity is genuine,
not a failure. Throughput: 50 volumes/second at $32^3$. The cross-domain
mapping: subbands are primes, subband energy is $p^{-\sigma}$, energy decay
across levels is $f(b)$, and anisotropy per level is
$\cos(\ln(p)\cdot\gamma)$.

\subsection{Turbulence Cascade Classification}

Multi-level 1D Haar DWT decomposes velocity signals into per-scale energy
and kurtosis profiles. The engine generates 29~features: cascade energy
ratios, K41 deviation (theoretical $2^{5/3}$ ratio), kurtosis growth rate
across scales, spectral flatness, cross-scale magnitude correlation, and
self-similarity. Validated on 5~synthetic regimes: laminar (sinusoidal),
Kolmogorov ($k^{-5/3}$ filtered noise), intermittent (Kolmogorov $+$ $5\times$
bursts), white noise (broadband), and tonal (machine-like harmonics). The
classifier achieves 5/5 correct using kurtosis sign (sub-Gaussian for
sinusoids), burst detection (extreme kurtosis for intermittency), and
spectral slope (negative for broadband). Limitation: Haar has poor frequency
resolution; the K41 cascade ratio is not recovered cleanly. Real turbulence
analysis requires Daubechies or Morlet wavelets. Throughput: 476
signals/second at 4096~samples. The cross-domain mapping: wavelet scales are
primes, $E(j)$ is $p^{-\sigma}$, cascade ratio is $f(b)$, and kurtosis per
scale is $\cos(\ln(p)\cdot\gamma)$.

\subsection{Image Spectral Analysis}

2D FFT with azimuthal averaging yields a radial power spectrum $P(f)$ from
grayscale images. The engine generates 29~features: spectral slope (natural
images follow $1/f^2$, per Field 1987), band energies, peak detection,
angular isotropy, spectral centroid, and forgery anomaly scoring. A forgery
detection module compares local vs.\ global spectral slope across
$16\times 16$ blocks to identify pasted regions. Validated on 6~synthetic
images: white noise (slope $\approx 0$), natural-like (slope $\approx -2$),
periodic texture, step edge, noisy natural, and forged (natural $+$ pasted
periodic block). Classifier: 6/6 correct. The forged image is correctly
identified as TEXTURED at the global level because the pasted periodic block
injects spectral peaks. Limitation: $64\times 64$ pixels; real forensics
needs megapixel images with multi-scale block analysis. Throughput: 816
images/second. The cross-domain mapping: spatial frequency is~$t$, spectral
peaks are zeros, $1/f^2$ decay is $f(b)$, and angular sector distribution is
$\cos(\ln(p)\cdot\gamma)$.

\subsection{Spectral Pathfinding}

Building on \texttt{graph\_engine.py} (\S7.9), this engine adds
pathfinding-specific features: spectral coordinates from Laplacian
eigenvectors, effective resistance via pseudoinverse, BFS vs.\ spectral
distance rank correlation, and algorithm recommendation. The engine generates
29~features. Validated on 5~graph types: $10\times 10$ grid
(A\_STAR\_SPECTRAL, rank correlation $0.71$), Erd\H{o}s--R\'enyi 100-node
(BFS, well-connected), bottleneck SBM (BRIDGE\_FIRST, high eigengap),
Barab\'asi--Albert (ROUTE\_HUBS, scale-free), and 50-node path
(A\_STAR\_SPECTRAL). The spectral heuristic is most valuable for graphs with
bottleneck structure where eigenvector embedding reveals global geometry that
BFS cannot see. For regular/expander graphs, BFS is already optimal.
Throughput: 22 graphs/second at $N=100$. The cross-domain mapping: node
index is~$t$, critical nodes (Fiedler) are zeros, spectral coordinates are
the torus embedding, and effective resistance is $f(b)$.

\subsection{Elliptic Curve Cryptographic Auditing}

For elliptic curves $y^2 = x^3 + ax + b$ over $\mathbb{F}_p$, brute-force
point counting gives $\#E(\mathbb{F}_p)$ and $a_p = p + 1 - \#E$. The
engine generates 29~features from group order factorization (Canal~A: largest
prime factor ratio, smoothness, cofactor), embedding degree (Cross:
MOV/Frey--R\"uck vulnerability), Frobenius eigenvalue structure, anomalous
flag ($\#E = p$), and a composite security score. Validated on 5~curves:
secure ($y^2 = x^3 + 3x + 2$ over $\mathbb{F}_{127}$, prime order~139,
$\text{emb}_k = 69$), Pohlig--Hellman (16-smooth order), anomalous
($\#E = p$), MOV-weak (supersingular, $\text{emb}_k = 2$), and
supersingular ($a_p = 0$). Classifier: 5/5 correct, with
anomalous/MOV/supersingular correctly flagged as WEAK with multiple issues.
Limitation: brute-force counting is $O(p^2)$, limiting~$p$ to ${\sim}500$.
Real cryptographic curves use $p > 2^{200}$. The cross-domain mapping: group
order is the spectral gap, factorization structure is Canal~A, embedding
degree is the Ramanujan bound, and $a_p/(2\sqrt{p})$ is
$\cos(\ln(p)\cdot\gamma)$.


% ══════════════════════════════════════════════════════════════════════
\section{Open Questions}
\label{sec:open}
% ══════════════════════════════════════════════════════════════════════

Note: \S\S7.9--7.11 were validated after the initial draft, confirming
zero-change transfer.

\begin{enumerate}[label=(\alph*)]
\item Does the $R^2$ of the factored form
  $\theta^* = f(b)\cdot\cos(\omega)$ improve with the Weil quadratic form as
  loss instead of the log-derivative? If $R^2 \to 1$, the decomposition is
  exact and $f(b)\cdot\cos(\omega)$ is the prolate eigenvector in closed
  form.

\item Is the third axis ($\cos(\ln(p)\cdot\gamma)$) sufficient, or does a
  fourth axis emerge at higher~$K$? The residual after $A+B+C$ projection is
  $55\%$ of norm at $K=30$. A neural network with more capacity might reduce
  this, or it might be irreducible noise.

\item Does the $V(\sigma)$ profile at $s = 1$ discriminate elliptic curve
  rank at scale? The toy (4~curves, $K=20$) is suggestive but far from
  conclusive. LMFDB provides thousands of curves with certified rank.

\item Can the harmonicity of $V(\sigma,t)$ be improved by using the Weil
  form instead of the log-derivative? The $0.07\%$ residual may be inherent
  to the partial product or may be reducible with better operators.
\end{enumerate}


% ══════════════════════════════════════════════════════════════════════
\section{Deliverables}
\label{sec:deliverables}
% ══════════════════════════════════════════════════════════════════════

\begin{enumerate}[label=(\alph*)]
\item \texttt{geometric\_engine.py} --- self-contained Python engine
  (numpy/scipy only), 8~modules, 29~features, ${\sim}8\,$ms per query.
  Supports $\zeta(s)$, $L(s,\chi)$, $L(s,E)$ with \texttt{EllipticCurveL}
  class. AGPL-3.0.
\item \texttt{sparam\_engine.py} --- S-parameter spectral engine (1709/s).
\item \texttt{motor\_diagnosis.py} --- vibration fault diagnosis (1132/s).
\item \texttt{graph\_engine.py} --- general graph spectral analysis
  (833/s at $N=100$).
\item \texttt{ast\_engine.py} --- Python AST spectral analysis,
  $29 + 15 = 44$ features (139/s).
\item \texttt{cavity\_engine.py} --- physical cavity optimization,
  8~shapes, application scoring (20/s).
\item \texttt{phonon\_engine.py} --- 1D chain phonon spectra, 12~materials,
  thermoelectric screening (478/s).
\item \texttt{colorimetry\_engine.py} --- visible spectrum analysis, metamer
  detection, $2\times$ XYZ discrimination (1677/s).
\item \texttt{wavelet\_3d\_engine.py} --- 3D Haar DWT, 6~synthetic volumes
  at $32^3$, structural classification (50/s).
\item \texttt{turbulence\_engine.py} --- 1D Haar DWT, 5~turbulence regimes,
  cascade analysis (476/s).
\item \texttt{image\_spectral\_engine.py} --- 2D FFT radial spectrum,
  forgery detection, 6~image types (816/s).
\item \texttt{pathfinding\_engine.py} --- spectral graph pathfinding,
  algorithm recommendation, builds on \texttt{graph\_engine} (22/s at
  $N=100$).
\item \texttt{crypto\_curve\_engine.py} --- elliptic curve security audit
  over $\mathbb{F}_p$, 5~vulnerability classes (274/s at $p=101$).
\end{enumerate}

All engines share the same 29-feature core architecture.

Additional deliverables:
\begin{enumerate}[resume,label=(\roman*)]
\item \texttt{envelope\_pipeline.py} (v1) and \texttt{envelope\_v2.py} ---
  optimization pipelines for envelope weights. L-BFGS and neural network
  variants.
\item 17 interactive HTML visualizations (Three.js r128): toroidal
  embeddings, torus modes, operators, landscapes, spectra.
\item \texttt{dual\_role\_f\_v3.docx} --- technical report on the three
  scales, $f(b)$, and Poincar\'e compactification, with all 7~overclaims
  documented and corrected.
\item \texttt{class\_number\_report.docx} --- the identity
  $S - k/2 = -h(-p)/2$, route of discovery, and 961-prime verification.
\end{enumerate}


% ══════════════════════════════════════════════════════════════════════
\section{Conclusion}
\label{sec:conclusion}
% ══════════════════════════════════════════════════════════════════════

The engine is the product; the results are the byproducts. No individual
result survived as mathematical novelty against the existing literature. What
survived is a methodology that eliminates impossible approaches at visual
speed, a decomposition that separates infrastructure from signal in any
spectrum, and a modular tool that generates geometric parameters without
statistical assumptions. The architecture transfers with zero code changes
from $\zeta(s)$ to elliptic curves, Cayley graphs, differential
S-parameters, vibration spectra, physical cavities, and visible-light
spectroscopy---because the spectrum is the spectrum.

The central finding, confirmed by three independent tests across two
$L$-function families, is that the arithmetic structure of $p-1$ is
orthogonal to the spectral structure of the Euler product phases. Shared
algebra does not imply shared information. The operator
$f(b) = b/(b+1)$ acts identically in both channels but transports
information in only one. This orthogonality constrains which mathematical
structures can and cannot contribute to the detection of zeta zeros, and may
inform the design of future operators.

The engine does not prove theorems, replace domain physics, predict temporal
evolution, or handle strongly nonlinear systems without wavelet
linearization. It generates features, not demonstrations. Its value is in
filtering: generating cheap geometric coordinates that concentrate expensive
computation on the few candidates that matter.

The phrase that closes the arc: $t$ parametrises the critical line.
Everything else is scaffolding.
\begin{equation*}
  p^{-s} = p^{-\sigma} \cdot e^{-i\ln(p)\cdot t}
\end{equation*}
This equation is a fact, not a choice. $p^{-\sigma}$ is the amplitude.
$e^{-i\omega t}$ is the phase. The zeros are the values of~$t$ where all
phases conspire. $f(b)$ organises the amplitude. $\cos(\omega)$ organises
the phase. Canal~A is orthogonal to both. Everything in this paper---the
embedding, the three scales, the potential, the 16~domains, the motor---is
scaffolding around this single factorisation.

A note on democratisation. The engine lowers the barrier of access: anyone
with Python and a browser can explore~$\zeta$, diagnose a turbine, or screen
a material. What it does not lower is the barrier of judgement. The
7~overclaims in~\S0 were produced by the same tool that produced the
surviving results. Detecting them required domain knowledge the tool does not
have. Accessible tools without accessible scepticism produce false
confidence. The real democratisation in this paper is not the motor---it
is~\S0. The errata-first protocol teaches that the tool fails, when it
fails, and how to detect the failure. \emph{Democratising scepticism is
worth more than democratising computation.}


% ══════════════════════════════════════════════════════════════════════
% BIBLIOGRAPHY
% ══════════════════════════════════════════════════════════════════════
\begin{thebibliography}{9}

\bibitem{Connes2026}
A.~Connes,
\textit{Spectral truncations of the Bost--Connes system and the prolate
operator},
arXiv:2602.04022, February 2026.

\bibitem{Conrad2005}
K.~Conrad,
\textit{Partial Euler products on the critical line},
Canadian J.\ Math.\ \textbf{57} (2005), no.~2, 267--297.

\bibitem{GuthMaynard2024}
L.~Guth and J.~Maynard,
\textit{New large value estimates for Dirichlet polynomials},
2024, accepted by Ann.\ of Math.\ (2025).

\bibitem{WaveFormer2025}
Al~Hasan, M.\,M., Zaman, M., Jawad, A., Santamaria-Pang, A.,
Lee, H.\,H., Tarapov, I., See, K., Imran, M.\,S., Roy, A.,
Pourmohammadi~Fallah, Y., Asadizanjani, N., and Forghani, R.
WaveFormer: A 3D Transformer with Wavelet-Driven Feature
Representation for Efficient Medical Image Segmentation.
In \textit{MICCAI 2025}, LNCS 15963, Springer, 2026.
\texttt{arXiv:2503.23764}.

\bibitem{WDM2024}
Friedrich, P., Wolleb, J., Bieder, F., Durrer, A., and Cattin, P.\,C.
WDM: 3D Wavelet Diffusion Models for High-Resolution Medical
Image Synthesis.
In \textit{DGM4MICCAI 2024}, LNCS 15224, pp.\ 11--21, Springer, 2025.
\texttt{arXiv:2402.19043}.
\texttt{doi:10.1007/978-3-031-72744-3\_2}.

\end{thebibliography}

\end{document}
